\bigskip\noindent{\Large\bfseries\color{apren} Test 2. Generadors, independència i bases}\par\medskip
\iniciTest{espvecT2}{c,a,d,b,c,a,b,d,c,b}
\begin{enumerate}
\pregunta Una combinació lineal de $v_1,\ldots,v_n$ és una expressió de la forma:
\begin{enumerate}
\opcio{a}{$v_1v_2\cdots v_n$}
\opcio{b}{$v_1+\cdots+v_n+n$}
\opcio{c}{$\lambda_1v_1+\cdots+\lambda_nv_n$}
\opcio{d}{$\lambda_1+\cdots+\lambda_n$}
\end{enumerate}\feedback
\pregunta El subespai generat per una família de vectors és:
\begin{enumerate}
\opcio{a}{El conjunt de totes les seves combinacions lineals}
\opcio{b}{El conjunt dels vectors de la família només}
\opcio{c}{Sempre tot l’espai}
\opcio{d}{Sempre el subespai nul}
\end{enumerate}\feedback
\pregunta Una família és linealment independent si:
\begin{enumerate}
\opcio{a}{Tots els vectors tenen norma 1}
\opcio{b}{Cap vector és zero i prou}
\opcio{c}{Genera tot l’espai}
\opcio{d}{L’única combinació lineal nul·la és la trivial}
\end{enumerate}\feedback
\pregunta Si una família conté el vector zero, aleshores és:
\begin{enumerate}
\opcio{a}{Sempre generadora}
\opcio{b}{Linealment dependent}
\opcio{c}{Linealment independent}
\opcio{d}{Una base}
\end{enumerate}\feedback
\pregunta Una base d’un espai vectorial és una família:
\begin{enumerate}
\opcio{a}{Només independent}
\opcio{b}{Només generadora}
\opcio{c}{Generadora i linealment independent}
\opcio{d}{Amb vectors ortogonals}
\end{enumerate}\feedback
\pregunta La dimensió d’un espai vectorial finit és:
\begin{enumerate}
\opcio{a}{El nombre de vectors de qualsevol base}
\opcio{b}{El nombre de tots els vectors de l’espai}
\opcio{c}{Sempre el nombre d’equacions d’un sistema}
\opcio{d}{La suma de les coordenades d’una base}
\end{enumerate}\feedback
\pregunta A $\mathbb R^3$, tres vectors linealment independents:
\begin{enumerate}
\opcio{a}{No poden generar $\mathbb R^3$}
\opcio{b}{Formen una base de $\mathbb R^3$}
\opcio{c}{Són necessàriament iguals}
\opcio{d}{Tenen suma nul·la}
\end{enumerate}\feedback
\pregunta Què indica que $v$ pertanyi a $\langle v_1,\ldots,v_n\rangle$?
\begin{enumerate}
\opcio{a}{$v$ és perpendicular a tots els $v_i$}
\opcio{b}{$v$ és un dels $v_i$}
\opcio{c}{$v=0$}
\opcio{d}{$v$ és combinació lineal dels $v_i$}
\end{enumerate}\feedback
\pregunta En un espai de dimensió $n$, una família de més de $n$ vectors és sempre:
\begin{enumerate}
\opcio{a}{Una base}
\opcio{b}{Independent}
\opcio{c}{Dependent}
\opcio{d}{Ortonormal}
\end{enumerate}\feedback
\pregunta Les coordenades d’un vector respecte d’una base són:
\begin{enumerate}
\opcio{a}{Sempre úniques només si el vector és zero}
\opcio{b}{Úniques per a cada vector}
\opcio{c}{No existeixen si la dimensió és finita}
\opcio{d}{Sempre totes iguals}
\end{enumerate}\feedback
\end{enumerate}